% Chapter 3: Methodology — shared infrastructure

\chapter{Methodology}

This thesis builds a shared Python infrastructure (\texttt{satellite-paraguay})
that supports six peer-reviewed papers. This chapter describes the shared
components.

\section{Repository Structure}

The Python package is organized as follows:

\begin{itemize}
    \item \texttt{src/satellite\_io} — Sentinel-2, Landsat, Planet download
    \item \texttt{src/paraguay\_admin} — 18 departamentos + 7,912 tiles
    \item \texttt{src/foundation\_models} — Prithvi, AlphaEarth, DINOv2
    \item \texttt{src/parcel\_analysis} — Catastro intersection + indigenous conflict
    \item \texttt{src/timeseries} — Multi-temporal stacking + change detection
    \item \texttt{src/evaluation} — F1, IoU, mAP, MAE, R$^2$ metrics
    \item \texttt{src/papers/} — 6 paper-specific pipelines
    \item \texttt{api/main.py} — FastAPI endpoint
    \item \texttt{dashboard/app.py} — Streamlit dashboard
\end{itemize}

\section{Data Sources}

The package integrates 14+ datasets (see \texttt{data/datasheets/}):

\begin{itemize}
    \item \textbf{Local:} paraguay-geodata (549 MB) — Catastro, OSM, climate, indigenous
    \item \textbf{Sentinel-2:} ESA Copernicus, 10m, 5-day
    \item \textbf{Landsat:} NASA, 30m, 16-day
    \item \textbf{MapBiomas Paraguay:} 30m, annual, 1985-2024
    \item \textbf{Hansen GFC:} 30m, annual, 2000-2023
    \item \textbf{OpenAQ:} Hourly air quality, global
    \item \textbf{Sentinel-5P:} Atmospheric NO$_2$, O$_3$, SO$_2$, CO
    \item \textbf{NASA FIRMS:} Real-time fire alerts
    \item \textbf{Verra VCS:} Carbon credit registry
    \item \textbf{GBIF:} Biodiversity observations
\end{itemize}

\section{Foundation Models}

We use pre-trained foundation models that are fine-tuned on Paraguay data:

\begin{itemize}
    \item \textbf{Prithvi (IBM-NASA):} Vision Transformer pre-trained on HLS data \citep{jakubik2023foundation}
    \item \textbf{AlphaEarth (Google DeepMind):} Earth observation embeddings \citep{lamahewage2026alphaearth}
    \item \textbf{DINOv2 (Meta):} Self-supervised vision features \citep{oquab2024dinov2}
    \item \textbf{LLaVA-1.6:} Vision-language model for indigenous place names \citep{liu2024llavanext}
    \item \textbf{YOLOv8:} Object detection for wildlife poaching \citep{jocher2023yolov8}
    \item \textbf{TimesFM:} Time-series foundation model \citep{timesfm2024}
    \item \textbf{Whisper:} Speech recognition for clinical scribe \citep{whisper2022}
\end{itemize}

\section{Baselines}

Each paper includes baseline comparisons:

\begin{itemize}
    \item P0011 Yvutu: Random Forest, U-Net from scratch, persistence
    \item P0100 Yvyra: Linear regression, Random Forest, persistence
    \item P0025 Yrupe: Persistence, mean, simple regression
    \item P0035 Tatakua: Mean forecast, persistence, linear trend
    \item P0026 Kai: YOLOv8n pretrained (no fine-tune)
    \item P0012 Yvy: Catastro-only baseline (no AI)
\end{itemize}

\section{Evaluation Framework}

We use standard metrics from the literature:

\begin{itemize}
    \item \textbf{Segmentation:} F1 macro, mean IoU
    \item \textbf{Detection:} mAP@0.5, mAP@0.5:0.95
    \item \textbf{Regression:} R$^2$, MAE, RMSE
    \item \textbf{Classification:} Accuracy, F1, AUC
\end{itemize}

\section{Reproducibility}

We ensure reproducibility via:

\begin{itemize}
    \item Random seeds for all libraries (\texttt{set\_seed(42)})
    \item Environment capture (\texttt{capture\_environment})
    \item Git commit hash + branch tracking
    \item MLflow experiment tracking
    \item DVC data versioning
\end{itemize}

\section{Ethical Considerations}

The package follows ethical principles:

\begin{itemize}
    \item \textbf{CARE Principles} for indigenous data (\texttt{C}, \texttt{A}, \texttt{R}, \texttt{E})
    \item \textbf{FAIR} data principles
    \item \textbf{Open-source licensing} (Apache 2.0, MIT, GPL-3.0)
    \item \textbf{Citation guidelines} for all sources
    \item \textbf{Risk register} with mitigations
\end{itemize}
